\documentclass[12pt]{article}
\usepackage[utf8]{inputenc}
\usepackage[T1]{fontenc}
\usepackage[margin=1in]{geometry}
\usepackage{setspace}
\setlength{\parindent}{2em}
\setlength{\parskip}{1.25em}
\renewcommand{\baselinestretch}{1.0}
\makeatletter
\newcommand{\@subtitle}{}
\newcommand{\subtitle}[1]{\renewcommand{\@subtitle}{#1}}
\renewcommand{\@maketitle}{%
  \newpage
  \null
  \vskip 1em%
  \begin{center}%
    {\LARGE \@title \par}%
    \ifx\@subtitle\@empty\else
      \vskip 0.5em%
      {\large \@subtitle \par}%
    \fi
  \end{center}%
  \par
  \vskip 1em}
\makeatother
\usepackage{xcolor}
\usepackage[colorlinks=true,urlcolor=blue!55!black,linkcolor=black]{hyperref}

\title{\textit{City of Boerne v. Flores}, 521 U.S. 507 (1997)}
\subtitle{Case Brief}
\author{}
\date{}

\begin{document}
\maketitle

\section{Memory Jogger}

In \href{https://www.courtlistener.com/opinion/118140/city-of-boerne-v-flores/}{\textit{City of Boerne v. Flores}}, the Supreme Court struck down the Religious Freedom Restoration Act of 1993 as applied to the states, holding that Congress's enforcement power under Section 5 of the Fourteenth Amendment is strictly remedial and preventive---requiring ``congruence and proportionality'' between the constitutional injury and the legislative means---rather than a power to redefine substantive constitutional rights.

\section{Detailed Case Facts}

St. Peter Catholic Church, built in 1923 in mission-revival style on a hill in Boerne, Texas (about 28 miles northwest of San Antonio), seated roughly 230 worshippers---too few for its growing parish, leaving 40 to 60 parishioners unable to attend some Sunday masses. In 1993, Archbishop Patrick F. Flores of San Antonio authorized the parish to plan an expansion of the church building. A few months later, the Boerne City Council enacted a historic preservation ordinance authorizing the city's Historic Landmark Commission to designate historic landmarks and districts and requiring Commission preapproval for any construction affecting designated properties. When Archbishop Flores applied for a building permit to enlarge St. Peter Catholic Church, city authorities denied the application on the ground that the church lay within a newly designated historic district.

Archbishop Flores challenged the permit denial under the Religious Freedom Restoration Act of 1993 (RFRA), 42 U.S.C. \S~2000bb \textit{et seq.} Congress had enacted RFRA by overwhelming bipartisan majorities in direct response to \textit{Employment Division, Department of Human Resources of Oregon v. Smith}, 494 U.S. 872 (1990), where the Supreme Court held that the First Amendment's Free Exercise Clause does not exempt religious objectors from neutral, generally applicable laws, declining to apply the strict-scrutiny balancing test of \textit{Sherbert v. Verner}, 374 U.S. 398 (1963). Explicitly aiming to ``restore the compelling interest test'' of \textit{Sherbert} and \textit{Wisconsin v. Yoder}, 406 U.S. 205 (1972), RFRA prohibited ``[g]overnment''---defined to include all federal, state, and local branches, agencies, and officials---from ``substantially burden[ing]'' a person's exercise of religion, even through a rule of general applicability, unless the government demonstrated that the burden (1) furthered a ``compelling governmental interest'' and (2) was the ``least restrictive means'' of furthering that interest (42 U.S.C. \S~2000bb-1).

To apply RFRA to state and local governments, Congress invoked Section 5 of the Fourteenth Amendment, which empowers Congress ``to enforce, by appropriate legislation, the provisions of this article''---including the Due Process Clause of Section 1, through which the Free Exercise Clause applies to the states under \textit{Cantwell v. Connecticut}, 310 U.S. 296 (1940). See \href{https://www.fjc.gov/history/cases/cases-that-shaped-the-federal-courts/city-boerne-v-flores}{Federal Judicial Center Summary} and \href{https://www.oyez.org/cases/1996/95-2074}{Oyez Case Record}.

\section{Procedural History}

Archbishop Flores filed suit against the City of Boerne in the U.S. District Court for the Western District of Texas, asserting multiple claims including a violation of RFRA. In \textit{Flores v. City of Boerne}, 877 F. Supp. 355 (W.D. Tex. 1995), the District Court held that Congress had exceeded its authority under Section 5 of the Fourteenth Amendment by enacting RFRA to overturn \textit{Smith}, concluding that the statute violated separation-of-powers principles.

The District Court certified its order for interlocutory appeal. The U.S. Court of Appeals for the Fifth Circuit reversed in \textit{Flores v. City of Boerne}, 73 F.3d 1352 (5th Cir. 1996), upholding RFRA as a valid exercise of Congress's Section 5 power to provide broader statutory protection for religious liberty than the constitutional floor recognized in \textit{Smith}. The City of Boerne petitioned for certiorari, which the Supreme Court granted at 519 U.S. 926 (1996). Following oral argument on February 19, 1997, the Supreme Court reversed the Fifth Circuit on June 25, 1997.

\section{Judicial Votes}

\begin{itemize}
  \item \textbf{Majority (6--3):} Justice Kennedy delivered the opinion of the Court, joined in full by Chief Justice Rehnquist and Justices Stevens, Thomas, and Ginsburg, and joined as to all but Part III-A-1 (the drafting history of the Fourteenth Amendment) by Justice Scalia.
  \item \textbf{Concurrences:}
  \begin{itemize}
    \item Justice Stevens filed a concurring opinion.
    \item Justice Scalia filed an opinion concurring in part, joined by Justice Stevens.
  \end{itemize}
  \item \textbf{Dissents:}
  \begin{itemize}
    \item Justice O'Connor filed a dissenting opinion, joined by Justice Breyer except as to the first paragraph of Part I.
    \item Justice Souter filed a dissenting opinion.
    \item Justice Breyer filed a dissenting opinion.
  \end{itemize}
\end{itemize}

\section{Holding}

\textbf{Judgment:} Reversed (6--3). The Religious Freedom Restoration Act of 1993 exceeds Congress's enforcement power under Section 5 of the Fourteenth Amendment and is unconstitutional as applied to state and local governments.

\textbf{Rules Announced:}
\begin{enumerate}
  \item Congress's Section 5 power ``to enforce'' the Fourteenth Amendment is strictly remedial and preventive, not substantive; Congress may not alter or expand the substantive meaning of constitutional rights as interpreted by the Supreme Court.
  \item To qualify as valid remedial or preventive legislation under Section 5, ``[t]here must be a congruence and proportionality between the injury to be prevented or remedied and the means adopted to that end.''
  \item RFRA fails the congruence-and-proportionality test because its legislative record lacks evidence of widespread, modern state laws enacted out of religious bigotry, and its sweeping, perpetual, nationwide strict-scrutiny mandate intrudes heavily into traditional state regulatory authority far out of proportion to any pattern of unconstitutional conduct under \textit{Smith}.
\end{enumerate}

\textbf{Holding vs. Dicta:} The invalidation of RFRA as applied to state and local governments under Section 5, along with the ``congruence and proportionality'' framework, is binding holding. Notably, \textit{Boerne} did not invalidate RFRA as applied to the federal government itself (which rests on Congress's Article I powers over federal laws, later confirmed in \textit{Gonzales v. O Centro Esp\'{i}rita Beneficente Uni\~{a}o do Vegetal}, 546 U.S. 418 (2006)), nor did the Court hold that Section 5 legislation must always contain geographic limits or sunset dates---noting only that such tailoring helps demonstrate proportionality when Congress prohibits otherwise constitutional state conduct.

\section{Analysis of Opinions}

\subsection{Majority Opinion (Justice Kennedy)}

Writing for the Court (\href{https://www.law.cornell.edu/supct/html/95-2074.ZO.html}{Cornell LII Full Text}), Justice Kennedy anchored the analysis in enumerated powers and judicial supremacy under \textit{Marbury v. Madison}, 5 U.S. (1 Cranch) 137 (1803). While acknowledging that Section 5 is a ``positive grant of legislative power'' (\textit{Katzenbach v. Morgan}, 384 U.S. 641, 651 (1966)) that allows Congress to enact preventive rules deterring unconstitutional conduct even if those rules prohibit some otherwise constitutional state actions, Kennedy insisted that the text of Section 5 grants only the power ``to enforce,'' not the power to determine what constitutes a constitutional violation. To distinguish permissible remedial measures from impermissible substantive redefinitions of constitutional rights, the Court announced its governing standard: ``There must be a congruence and proportionality between the injury to be prevented or remedied and the means adopted to that end.''

In Part III-A-1 (which Justice Scalia did not join), Kennedy surveyed the drafting history of the Fourteenth Amendment in the 39th Congress. He highlighted the House's February 1866 rejection of Representative John Bingham's initial draft---which would have empowered Congress to make all necessary and proper laws to secure privileges, immunities, and equal protection---because lawmakers feared it would centralize plenary legislative authority in Congress at the expense of state sovereignty and judicial interpretation. The revised Amendment ratified in 1868 instead enacted self-executing limits on the states in Section 1 and a corrective, remedial enforcement power in Section 5. Turning to precedent in Parts III-A-2 and III-A-3, Kennedy reconciled the \textit{Civil Rights Cases}, 109 U.S. 3 (1883), \textit{South Carolina v. Katzenbach}, 383 U.S. 301 (1966), \textit{Oregon v. Mitchell}, 400 U.S. 112 (1970), and \textit{Katzenbach v. Morgan}, rejecting the ``ratchet theory'' reading of \textit{Morgan} on the ground that allowing Congress to redefine the Fourteenth Amendment would place the Constitution ``on a level with ordinary legislative acts'' and circumvent Article V.

Applying the congruence-and-proportionality test in Part III-B, Kennedy contrasted RFRA with the Voting Rights Act of 1965. Whereas the Voting Rights Act responded to a century of documented, pervasive, intentional racial discrimination in voting and used targeted geographic coverage and bailout/termination mechanisms, RFRA's legislative record contained no examples in the prior 40 years of generally applicable laws passed out of religious bigotry---only anecdotal accounts of incidental burdens from neutral laws such as autopsy rules and zoning ordinances. Moreover, RFRA's means were grossly disproportionate: it applied to every federal, state, and local agency and law (past and future), contained no sunset date or geographic limit, and imposed ``the most demanding test known to constitutional law'' (including a least-restrictive-means requirement not even used in pre-\textit{Smith} case law). Concluding that RFRA contradicted both separation of powers and the federal-state balance, the Court reversed.

\subsection{Concurring Opinion (Justice Stevens)}

In a brief separate concurrence (\href{https://supreme.justia.com/cases/federal/us/521/507/}{Justia Full Text}), Justice Stevens argued that RFRA independently violated the First Amendment's Establishment Clause. If the historic landmark in Boerne had been an atheist-owned museum or art gallery, it would have had no federal statutory claim to an exemption from the preservation ordinance. Because RFRA granted religious institutions ``a legal weapon that no atheist or agnostic can obtain,'' Stevens concluded that the statute constituted an impermissible governmental preference for religion over irreligion.

\subsection{Concurring Opinion in Part (Justice Scalia)}

Joined by Justice Stevens, Justice Scalia joined the Court's opinion except Part III-A-1 (consistent with his skepticism of legislative history) and wrote separately to refute Justice O'Connor's historical attack on \textit{Employment Division v. Smith}. Scalia argued that early colonial charters, state constitutions, and the Northwest Ordinance protected only against laws targeting religion ``on account of'' or ``in respect of'' faith, and that their ``peace and safety'' provisos meant obedience to general laws governing conduct---as reflected in \textit{Queen v. Lane} (1704) and Noah Webster's 1828 dictionary. He contended that early legislative accommodations (such as Quaker oath and conscription exemptions) showed only that legislatures \textit{may} grant exemptions, not that courts are constitutionally required to do so, and emphasized that not a single founding-era appellate decision had invalidated a neutral, generally applicable statute on free-exercise grounds. Ultimately, Scalia framed the issue as democratic self-governance: whether elected legislatures or courts should decide whether churches are exempt from zoning laws.

\subsection{Dissenting Opinion (Justice O'Connor)}

Joined by Justice Breyer except as to the first paragraph of Part I, Justice O'Connor expressly agreed with Part III-A of the majority's Section 5 analysis---including the congruence-and-proportionality standard and the principle that Congress cannot substantively redefine constitutional rights. However, she dissented because the majority measured RFRA against \textit{Smith}, which she argued was ``demonstrably wrong'' and should be overruled after reargument. Undertaking an extensive historical survey of colonial charters (Maryland 1649, Rhode Island 1663), revolutionary-era state constitutions, the Virginia Declaration of Rights debate between George Mason and James Madison, early accommodations for Quakers and other dissenters, and writings of Madison, Jefferson, Washington, and Ellsworth, O'Connor argued that the Free Exercise Clause was originally understood as an affirmative guarantee of religious conduct against unnecessary governmental interference, not merely an antidiscrimination rule.

\subsection{Dissenting Opinion (Justice Souter)}

Justice Souter expressed ``serious doubts'' about the precedential soundness of \textit{Smith}---doubts intensified by Justice O'Connor's historical analysis---and noted that the Court had never received full adversarial briefing on the merits of the \textit{Smith} rule. Because deciding RFRA's validity while uncritically assuming \textit{Smith}'s correctness left free-exercise doctrine in ``intolerable tension,'' Souter advocated setting the case for reargument or dismissing the writ of certiorari as improvidently granted.

\subsection{Dissenting Opinion (Justice Breyer)}

In a one-paragraph dissent, Justice Breyer agreed with Justice O'Connor that the Court should order reargument on whether \textit{Smith} was correctly decided. Unlike O'Connor, however, Breyer declined to endorse the majority's Section 5 analysis in the first paragraph of Part I, finding it unnecessary to decide whether Section 5 would authorize RFRA if \textit{Smith} were correct.

\section{Examples: Future Applications}

\begin{enumerate}
  \item \textbf{Same-Side (Invalidated): Abrogating State Sovereign Immunity Under the ADA for Non-Suspect Classifications.} Congress enacts a federal statute allowing private individuals to sue states for money damages for failing to provide workplace accommodations for disabilities, relying on Section 5 of the Fourteenth Amendment. Under \textit{Boerne}, because disability classifications receive only rational-basis review under the Equal Protection Clause, a broad statutory accommodation mandate unsupported by a documented pattern of irrational, unconstitutional state employment discrimination lacks congruence and proportionality and cannot validly abrogate state Eleventh Amendment immunity. (The Court applied \textit{Boerne} in precisely this way in \textit{Board of Trustees of the University of Alabama v. Garrett}, 531 U.S. 356 (2001).)
  \item \textbf{Same-Side (Invalidated): Federal Statute Redefining Constitutional Property Rights Against States.} Congress passes a nationwide statute prohibiting states and municipalities from enforcing any zoning or environmental regulation that reduces a landowner's property value by more than 15\% without paying compensation, invoking Section 5 to enforce the Fourteenth Amendment's Takings Clause. Under \textit{Boerne}, this statute exceeds Section 5 because it substantively rewrites the Court's regulatory-takings standard (\textit{Penn Central}) rather than remedying a documented pattern of unconstitutional state takings.
  \item \textbf{Opposite-Side (Upheld): Targeted Family-Leave Mandate Combating State Gender Stereotypes.} Congress enacts a narrowly tailored federal statute guaranteeing state employees up to 12 weeks of unpaid family-care leave to combat a well-documented, longstanding pattern of unconstitutional gender-based discrimination in state leave policies. Under \textit{Boerne}, because gender classifications trigger heightened scrutiny and Congress compiled an extensive record of unconstitutional state gender stereotyping, the narrowly focused leave remedy satisfies congruence and proportionality. (The Court upheld the FMLA's family-care provision on this exact rationale in \textit{Nevada Department of Human Resources v. Hibbs}, 538 U.S. 721 (2003).)
  \item \textbf{Opposite-Side (Upheld): Religious Land Use and Institutionalized Persons Act (RLUIPA) Under Spending and Commerce Powers.} Following \textit{Boerne}, Congress enacts a narrower statute imposing RFRA-like strict scrutiny on state prisons and local zoning boards only when the state program receives federal funds or when the land-use burden affects interstate commerce or involves individualized assessments. Because Congress relies on its Article I Spending and Commerce Clauses (and targets individualized assessments already subject to heightened scrutiny under \textit{Smith}) rather than a sweeping Section 5 mandate, the statute avoids \textit{Boerne}'s Section 5 bar. (The Court upheld RLUIPA's institutionalized-persons provision against an Establishment Clause challenge in \href{https://supreme.justia.com/cases/federal/us/544/709/}{\textit{Cutter v. Wilkinson}}, 544 U.S. 709 (2005).)
  \item \textbf{Fence-Sitter (Close Call): Reauthorizing Preclearance or Nationwide Voting-Access Rules Under the Fourteenth and Fifteenth Amendments.} Congress enacts a new federal voting-rights statute prohibiting states nationwide from adopting voter-ID or absentee-ballot restrictions that produce a statistically significant racial disparity, without geographic limitations or a sunset clause, citing recent hearings on minority turnout gaps. Under \textit{Boerne}, this is a close question: voting involves a fundamental right and racial discrimination triggers strict scrutiny (where Congress receives greater latitude under \textit{South Carolina v. Katzenbach}), yet the absence of geographic targeting, sunset provisions, or proof of widespread intentional unconstitutional discrimination invites a serious congruence-and-proportionality challenge.
\end{enumerate}

\section{Critique}

\textit{City of Boerne v. Flores} is a watershed federalism and separation-of-powers ruling that has attracted sharp, cross-ideological criticism from both progressive and originalist/textualist scholars, alongside strong support from judicial federalists.

From a \textbf{progressive and democratic-constitutionalist perspective}, scholars such as Robert C. Post and Reva B. Siegel (\textit{Protecting the Constitution from the People: Juricentric Restrictions on Section Five Power}, 78 Ind. L.J. 1 (2003); \textit{Legislative Constitutionalism and Section Five Power}, 112 Yale L.J. 1943 (2003)) criticize \textit{Boerne} for entrenching a rigid ``juricentric'' model of constitutional interpretation. They argue that the Court transformed \textit{Marbury v. Madison} into an aggressive doctrine of judicial exclusivity, treating Congress like an inferior administrative agency that must build a trial-like evidentiary record to justify civil rights legislation. Progressive critics further contend that \textit{Boerne}'s congruence-and-proportionality test proved doctrinally malleable and politically consequential, enabling the Rehnquist Court to strike down key remedial provisions of the Violence Against Women Act (\textit{United States v. Morrison}), the Age Discrimination in Employment Act (\textit{Kimel v. Florida Board of Regents}), and Title I of the Americans with Disabilities Act (\textit{Garrett}). See also \href{https://scholarlycommons.law.northwestern.edu/cgi/viewcontent.cgi?article=1349&context=nulr}{Northwestern University Law Review on Section 5 History}.

From an \textbf{originalist and textualist perspective}, prominent conservative and originalist scholars---most notably Michael W. McConnell in his landmark critique, \textit{Institutions and Interpretation: A Critique of City of Boerne v. Flores}, 111 Harv. L. Rev. 153 (1997), as well as Akhil Reed Amar and commentators cited by the \href{https://www.theusconstitution.org/wp-content/uploads/2017/12/Shield_of_National_Protection.pdf}{Constitutional Accountability Center}---challenge both the text and history of \textit{Boerne}. First, textually, Section 5 uses the phrase ``appropriate legislation,'' which since \textit{McCulloch v. Maryland} (1819) and \textit{Ex parte Virginia} (1880) signaled deferential means-ends rationality rather than the heightened ``congruence and proportionality'' scrutiny \textit{Boerne} invented. Second, historically, the Reconstruction Congress that drafted the Fourteenth Amendment deeply distrusted the Supreme Court after \textit{Dred Scott v. Sandford} and intended Congress---not solely the judiciary---to play a primary, active role in defining and securing Fourteenth Amendment liberties. McConnell also argues that Kennedy misread the 1866 rejection of the Bingham draft: lawmakers objected to granting Congress plenary power to displace state law absent state action, not to Congress exercising independent interpretive judgment when enforcing Section 1's guarantees.

Conversely, \textbf{federalism defenders}---such as Jeffrey S. Sutton (who argued as amicus for Ohio in \textit{Boerne}; see \href{https://fedsoc.org/case/city-of-boerne-v-flores}{Federalist Society Case Summary}) and William Van Alstyne---praise \textit{Boerne} as an indispensable guardrail for both enumerated powers and constitutional supremacy. If Congress could invoke Section 5 to overrule Supreme Court interpretations of constitutional rights whenever it disagreed with them, the federal government would acquire a general police power over state and local governance (from historic zoning to criminal law), and ordinary legislative majorities could effectively amend the Constitution without satisfying Article V. As subsequent legislation such as RLUIPA demonstrated (see \href{https://nyulawreview.org/wp-content/uploads/2018/08/NYULawReview-75-2-Proctor.pdf}{NYU Law Review on Spending Power Workarounds}), \textit{Boerne} forced Congress to rely on narrower, enumerated Article I powers rather than an unbounded Section 5 mandate.

\section{Key Quotations}

\begin{quote}
``Legislation which alters the meaning of the Free Exercise Clause cannot be said to be enforcing the Clause. Congress does not enforce a constitutional right by changing what the right is. It has been given the power `to enforce,' not the power to determine what constitutes a constitutional violation.'' --- Justice Kennedy, Opinion of the Court, 521 U.S. at 519.
\end{quote}

\begin{quote}
``There must be a congruence and proportionality between the injury to be prevented or remedied and the means adopted to that end. Lacking such a connection, legislation may become substantive in operation and effect.'' --- Justice Kennedy, Opinion of the Court, 521 U.S. at 520.
\end{quote}

\begin{quote}
``Shifting legislative majorities could change the Constitution and effectively circumvent the difficult and detailed amendment process contained in Article V.'' --- Justice Kennedy, Opinion of the Court, 521 U.S. at 529.
\end{quote}

\begin{quote}
``Whether the Church would actually prevail under the statute or not, the statute has provided the Church with a legal weapon that no atheist or agnostic can obtain. This governmental preference for religion, as opposed to irreligion, is forbidden by the First Amendment.'' --- Justice Stevens, concurring, 521 U.S. at 537.
\end{quote}

\begin{quote}
``Rather, the Clause is best understood as an affirmative guarantee of the right to participate in religious practices and conduct without impermissible governmental interference, even when such conduct conflicts with a neutral, generally applicable law.'' --- Justice O'Connor, dissenting, 521 U.S. at 546.
\end{quote}

\end{document}
